\chapter{Prêts, CLO et titrisation}\label{ch:m2:loans-clos-and-securitisation}

À la fin de 2025, les entreprises américaines devaient 1.55~billion de dollars de prêts institutionnels
à effet de levier, soit 9.2~\% de plus qu'un an plus tôt, selon le décompte de la Réserve fédérale.
Il s'agit de prêts accordés à des sociétés très endettées, et ces prêts ne sont pas conservés par les banques qui les organisent~:
ils sont vendus à des investisseurs institutionnels, et beaucoup finissent dans des véhicules qui achètent un
portefeuille diversifié de ces prêts et le financent en émettant des titres de créance de différents rangs de séniorité. Les documents de chaque véhicule précisent, au dollar près, qui est payé en premier à chaque
trimestre, et comportent des tests qui, lorsque les prêts commencent à faire défaut, coupent la trésorerie
versée aux investisseurs les plus subordonnés et la dirigent vers les plus prioritaires. Ce chapitre décrit les
prêts, la \omterm{def:m2:loans-clos-and-securitisation:sec}{titrisation}, la \omterm{def:m2:loans-clos-and-securitisation:clo}{collateralised loan obligation} et sa \omterm{def:m2:loans-clos-and-securitisation:clo}{cascade des paiements}, et qui
détient quelle part~; la construction porte sur la cascade elle-même.

\section{Le marché des prêts à effet de levier}

\begin{definition}[Prêt à effet de levier, term loan B, covenant-lite]\label{def:m2:loans-clos-and-securitisation:loan}
Un \emph{prêt à effet de levier}\index{prêt à effet de levier} est un prêt, en général senior et garanti sur
les actifs de l'emprunteur et rémunéré à taux variable, soit un taux de référence majoré d'un spread, consenti à une société
dont la dette est élevée au regard de ses bénéfices ou dont la notation se situe sous l'investment grade.
Un \emph{term loan B}\index{term loan B} est la tranche d'un tel prêt structurée pour être
vendue à des investisseurs institutionnels~: une maturité longue, un amortissement limité avant l'échéance, et
négociable. Un \emph{covenant-lite}\index{covenant-lite} est un prêt sans
covenants de maintien, c'est-à-dire des ratios financiers que l'emprunteur doit respecter à chaque date de contrôle~; il ne conserve que
des covenants d'incurrence, qui sont contrôlés lorsque l'emprunteur entreprend une action, par exemple
contracter un endettement supplémentaire.
\end{definition}

Une banque structure le prêt, par exemple pour financer le rachat d'une société par un fonds de
private equity, et la syndique~: elle en cède l'essentiel à des fonds de prêt, des CLO et d'autres
investisseurs, qui le négocient ensuite sur le marché secondaire. Le prêt étant à taux variable, la
charge d'intérêts de l'emprunteur suit le taux directeur du chapitre~1~; parce que le prêt est
senior et généralement garanti, ses prêteurs ont tendance à recouvrer davantage après un défaut que
les détenteurs d'obligations le font (les enchères du chapitre~23 le montrent)~; et lorsqu'il est \omterm{def:m2:loans-clos-and-securitisation:loan}{covenant-lite}, ils
disposent de moins d'alertes précoces avant qu'un défaut ne survienne.

La Réserve fédérale suit ce marché dans ses rapports sur la stabilité financière. En
mai 2026, elle a relevé que la part des nouveaux prêts accordés à des sociétés dont la dette représente au moins quatre fois
leurs bénéfices, avait augmenté et restait au-dessus de sa médiane historique, et que la
couverture des intérêts de l'emprunteur médian était proche de son plus bas historique, et que les défauts, y compris les
restructurations amiables, à savoir les prêts renégociés lorsque l'emprunteur était en difficulté, restaient relativement
élevés.

\begin{dated}{2026-09}{Les marchés des prêts à effet de levier et des CLO}\label{dat:m2:loans-clos-and-securitisation:size}
Encours des prêts institutionnels à effet de levier aux États-Unis~: 1\,549~milliards de dollars à la
fin de 2025, en hausse de 9.2~\% sur un an et de 12.6~\% par an en moyenne depuis 2000
(Federal Reserve, Financial Stability Report, mai 2026, à partir des données PitchBook LCD).
CLO~: 264~milliards de dollars en 2011, 617~milliards de dollars en 2018 (Federal Reserve, 2019, à partir des
données SIFMA)~; aucun chiffre public plus récent n'a été recueilli pour ce livre.
\end{dated}

\section{Titrisation}

\begin{definition}[Titrisation, entité ad hoc (SPV)]\label{def:m2:loans-clos-and-securitisation:sec}
La \emph{titrisation}\index{titrisation} est le regroupement de prêts ou d'autres
actifs financiers et la cession des créances sur leurs flux de trésorerie sous forme de titres. Une
\emph{entité ad hoc}\index{entité ad hoc} (special-purpose vehicle, SPV) est une société créée
uniquement pour détenir le portefeuille et émettre les titres, de sorte que celui-ci soit séparé de
l'entité qui a initié ou cédé les actifs, et que les créances des investisseurs ne dépendent que de lui.
\end{definition}

La \omterm{def:m2:loans-clos-and-securitisation:sec}{titrisation} transforme un portefeuille d'actifs en titres présentant des risques différents. Le
pass-through hypothécaire du chapitre 12 répartit les flux au prorata~; la plupart des autres
\omterm{def:m2:loans-clos-and-securitisation:sec}{titrisations} les découpent en tranches, comme le font les tranches d'indice du chapitre 24, de sorte que certains
investisseurs sont payés en premier et d'autres absorbent les premières pertes. Les cartes de crédit, les prêts automobiles,
les prêts hypothécaires commerciaux et les \omterm{def:m2:loans-clos-and-securitisation:loan}{prêts à effet de levier} sont tous titrisés. La Réserve fédérale
observe que la \omterm{def:m2:loans-clos-and-securitisation:sec}{titrisation} peut accroître l'effet de levier du système financier, parce que les
véhicules sont soumis à des règles moins contraignantes, comme l'obligation de rétention du risque, que les
exigences de fonds propres des banques.

La rétention du risque oblige le sponsor d'une \omterm{def:m2:loans-clos-and-securitisation:sec}{titrisation} à conserver une part de son risque de
crédit, soit 5~\% selon la règle américaine adoptée en vertu du Dodd--Frank Act, de sorte qu'il ait
un intérêt dans la qualité de ce qu'il vend. Pour une catégorie de CLO, elle ne s'applique pas~: en février 2018, la
cour d'appel fédérale de Washington a jugé que les gérants de CLO de marché ouvert, qui
n'initient les prêts ni ne les détiennent avant le véhicule, ne sont pas les
titriseurs visés par la règle, et celle-ci a été annulée en tant qu'elle leur était appliquée.

\section{La cascade des paiements du CLO et ses tests}

\begin{definition}[Collateralised loan obligation (CLO), cascade des paiements]\label{def:m2:loans-clos-and-securitisation:clo}
Un \emph{collateralised loan obligation}\index{collateralised loan obligation} (CLO)
est une \omterm{def:m2:loans-clos-and-securitisation:sec}{entité ad hoc} qui détient un portefeuille diversifié et géré activement de \omterm{def:m2:loans-clos-and-securitisation:loan}{prêts à effet de levier}
et le finance en émettant des titres notés, de rangs de séniorité décroissants, et des titres de capital non notés. Sa
\emph{cascade des paiements}\index{cascade des paiements} est l'ordre, fixé dans ses documents,
dans lequel les fonds encaissés au titre des prêts sont versés à chaque date de paiement.
\end{definition}

\begin{definition}[Test de surdimensionnement, test de couverture des intérêts, tranche equity]\label{def:m2:loans-clos-and-securitisation:tests}
Un \emph{test de surdimensionnement}\index{test de surdimensionnement} compare la
valeur au pair des prêts à celle des titres en circulation jusqu'à une classe donnée~; un
\emph{test de couverture des intérêts}\index{test de couverture des intérêts} compare les intérêts
encaissés avec ceux dus sur ces titres. Chacun dispose d'un ratio de déclenchement, et lorsqu'un
ratio passe sous ce seuil, les fonds qui auraient autrement poursuivi leur chemin dans la cascade servent à rembourser les
titres les plus prioritaires à la place. La \emph{tranche equity}\index{tranche equity}
d'une \omterm{def:m2:loans-clos-and-securitisation:sec}{titrisation} est sa créance la plus subordonnée, non notée, qui reçoit les fonds restants
une fois toutes les autres créances réglées.
\end{definition}

Un gérant sélectionne les prêts, les négocie et, pendant une période de réinvestissement de plusieurs
années, remplace ceux qui sont remboursés ou font défaut~; les titres versent un taux variable
majoré d'un spread fixé à l'émission, et les titres de capital reçoivent la différence entre ce que
rapportent les prêts et ce que coûtent les titres, soit l'\emph{excédent de spread}, tant que les tests sont
respectés (\cref{fig:m2:loans-clos-and-securitisation:structure}).

\begin{omfigure}
\centering
\begin{tikzpicture}[font=\small,
  box/.style={draw,thick,rounded corners=2pt,align=center,minimum height=0.9cm,inner sep=3pt},
  arr/.style={-{Stealth[length=2.5mm]},thick}]
\node[box,draw=gray,fill=gray!8,minimum width=3cm] (loans) at (0,0) {un portefeuille\\diversifié de prêts\\à effet de levier\\(500~millions\\de dollars au pair)};
\node[box,draw=omThm,fill=omThm!8,minimum width=2.6cm] (spv) at (5.4,0) {CLO (SPV)\\géré par un gérant};
\draw[arr] (loans) -- node[above,font=\footnotesize] {intérêts} node[below,font=\footnotesize] {et principal} (spv);
\foreach \n/\t/\y/\c in {aaa/{AAA 310 (62~\%)}/1.6/omDef, aa/{AA 55 (11~\%)}/0.8/omDef, a/{A 30 (6~\%)}/0/omProp, bbb/{BBB 30 (6~\%)}/-0.8/omProp, bb/{BB 25 (5~\%)}/-1.6/omProp}
  \node[box,draw=\c,fill=\c!8,minimum width=3.2cm,minimum height=0.6cm] (\n) at (10.2,\y) {\t};
\node[box,draw=gray,fill=gray!25,minimum width=3.2cm,minimum height=0.6cm] (eq) at (10.2,-2.4) {titres de capital 50 (10~\%)};
\foreach \n in {aaa,aa,a,bbb,bb,eq} \draw[arr] (spv.east) -- (\n.west);
\node[font=\footnotesize,align=left,anchor=west] at (12.25,1.2) {payée\\en premier,\\spread minimal};
\node[font=\footnotesize,align=left,anchor=west] at (12.25,-2.0) {payée\\en dernier,\\première perte};
\end{tikzpicture}

\omcaption{Le CLO illustratif de ce chapitre~: un portefeuille géré de \omterm{def:m2:loans-clos-and-securitisation:loan}{prêts à effet de levier} détenus
par une \omterm{def:m2:loans-clos-and-securitisation:sec}{entité ad hoc} (SPV) et financée par cinq classes d'obligations notées et 10~\%
titres de capital, en millions de dollars. Les flux de trésorerie descendent la structure dans l'ordre~; les pertes
sont absorbées à partir du bas. Schéma~; la structure est illustrative.}\label{fig:m2:loans-clos-and-securitisation:structure}
\end{omfigure}

Chaque trimestre, les intérêts encaissés couvrent d'abord la commission senior du gérant, puis les intérêts
de chaque classe à tour de rôle. Une fois une classe payée, ses tests sont vérifiés. Si un
test échoue, les intérêts restants servent à rembourser les obligations seniors jusqu'à ce que le
ratio soit revenu à son seuil, et seul ce qui reste alors descend dans la cascade
(\cref{fig:m2:loans-clos-and-securitisation:waterfall}). Les intérêts impayés sur les obligations
subordonnées ne constituent pas un défaut~: ils sont ajoutés à leur encours et versés plus tard s'il est possible.

\begin{omfigure}
\centering
\begin{tikzpicture}[font=\small,
  wstep/.style={draw,thick,rounded corners=2pt,align=center,minimum width=4.4cm,minimum height=0.62cm,inner sep=2pt},
  test/.style={draw,thick,diamond,aspect=2.6,align=center,inner sep=1pt,font=\footnotesize},
  arr/.style={-{Stealth[length=2.2mm]},thick}]
\node[wstep,draw=gray,fill=gray!8] (in) at (0,0) {intérêts encaissés};
\node[wstep,draw=gray,fill=gray!8] (fee) at (0,-1.0) {commission senior};
\node[wstep,draw=omDef,fill=omDef!8] (s) at (0,-2.0) {intérêts AAA, puis AA};
\node[test,draw=omThm] (t1) at (0,-3.25) {tests AA\\réussis~?};
\node[wstep,draw=omProp,fill=omProp!8] (m) at (0,-4.5) {intérêts A, BBB, BB,\\chacun suivi de son test};
\node[wstep,draw=gray,fill=gray!25] (eq) at (0,-5.7) {titres de capital~: ce qui reste};
\node[wstep,draw=omThm,fill=omThm!8,minimum width=3.6cm] (div) at (6,-3.25) {rembourser les obligations AAA jusqu'à ce que\\le test soit rétabli};
\draw[arr] (in) -- (fee); \draw[arr] (fee) -- (s); \draw[arr] (s) -- (t1);
\draw[arr] (t1) -- node[left,font=\footnotesize] {oui} (m); \draw[arr] (m) -- (eq);
\draw[arr] (t1) -- node[above,font=\footnotesize] {non} (div);
\draw[arr] (div.south) |- node[pos=0.25,right,font=\footnotesize,align=left] {reste, le cas échéant,\\se poursuit} (m.east);
\draw[arr,dashed,omThm] (m.west) -- ++(-1.4,0) |- node[pos=0.25,left,font=\footnotesize,align=right] {test junior\\échoue~:\\détournement} (s.west);
\end{tikzpicture}

\omcaption{La cascade des intérêts du CLO de ce chapitre. Après le paiement de chaque classe, ses
tests sont vérifiés~; un échec détourne les intérêts restants pour rembourser les obligations seniors
jusqu'à ce que le test soit rétabli, et seul ce qui reste descend jusqu'aux titres de capital.
Schéma.}\label{fig:m2:loans-clos-and-securitisation:waterfall}
\end{omfigure}

\begin{example}[Le spread excédentaire]\label{ex:m2:loans-clos-and-securitisation:base}
Le CLO du chapitre détient 500~millions de dollars de prêts rémunérés 3.5~\% au-dessus d'un taux
variable de 4~\%, soit 37.5~millions de dollars par an. La commission du gérant, de 0.45~\%, coûte 2.25~millions de dollars, et
les obligations, de l'AAA à 1.30~\% au-dessus jusqu'au BB à 6.00~\% au-dessus, coûtent 26.14~millions de dollars. Les
titres de capital, 50~millions de dollars, reçoivent 9.11~millions de dollars par an, soit 18.2~\% de leur investissement~;
sans défaut et avec des prêts vendus au pair au bout de cinq ans, leur taux de
rendement interne est de 19.5~\%. Les ratios de surdimensionnement partent de 1.370, 1.266, 1.176 et
1.111 aux niveaux AA, A, BBB et BB, face à des seuils de 1.25, 1.18, 1.10 et 1.05.
\end{example}

\begin{proposition}[Le coussin de surdimensionnement]\label{prop:m2:loans-clos-and-securitisation:cushion}
Si les obligations jusqu'à une classe donnée totalisent $D$, que les prêts ont un pair $P$ et que le
seuil de la classe est $\tau$, le test échoue dès que le pair a baissé de plus de $P - \tau D$. Avec
des recouvrements réinvestis au pair, chaque unité de pair défaillant réduit le pair du portefeuille de
$1 - R$, si bien que le test échoue dès que les défauts cumulés dépassent $(P - \tau D)/(1 - R)$.
\end{proposition}

\begin{proof}
Le test s'écrit $P'/D \geq \tau$. Avant tout détournement, $D$ reste inchangé, donc il échoue lorsque
$P' < \tau D$, soit une baisse de $P - \tau D$. Un défaut de $x$ retire $x$ de pair et le
recouvrement réinvesti restitue $Rx$, soit une baisse nette de $(1-R)x$.
\end{proof}

Pour le test BB, $P - \tau D = 500 - 1.05 \times 450 = 27.5$~: avec un recouvrement de 70~\%,
91.7~millions de dollars de défauts, soit 18.3~\% du portefeuille, le font échouer. Une fois qu'il échoue, le
rétablissement exige à peu près le pair perdu chaque trimestre, divisé par le seuil, sur les intérêts
qui seraient descendus dans la cascade~; lorsque cela dépasse ce qui reste après les
intérêts des obligations, les titres de capital ne reçoivent rien
(\cref{fig:m2:loans-clos-and-securitisation:paths}).

\begin{omfigure}
\begin{tikzpicture}
\begin{axis}[width=11cm,height=4.6cm,
  xlabel={trimestre},ylabel={ratio OC au niveau BB},
  tick label style={font=\small},label style={font=\small},
  xmin=0.5,xmax=20.5,ymin=1.0,ymax=1.12,grid=major,grid style={gray!25},
  yticklabel style={/pgf/number format/.cd,fixed,precision=2},
  legend columns=4,legend style={font=\footnotesize,at={(0.5,-0.32)},anchor=north,draw=none,fill=none,/tikz/every even column/.append style={column sep=8pt}},
  table/col sep=comma]
\addplot[omDef,very thick] table[x=quarter,y=oc0]{figdata/markets-2/25-loans-clos-and-securitisation/paths.csv};
\addplot[omThm,very thick] table[x=quarter,y=oc1]{figdata/markets-2/25-loans-clos-and-securitisation/paths.csv};
\addplot[omProp,very thick] table[x=quarter,y=oc2]{figdata/markets-2/25-loans-clos-and-securitisation/paths.csv};
\addplot[black,thick,dashed] coordinates {(0.5,1.05) (20.5,1.05)};
\legend{2~\% par an,4.8~\%,7~\%,seuil 1.05}
\end{axis}
\end{tikzpicture}

\begin{tikzpicture}
\begin{axis}[width=11cm,height=4.6cm,
  xlabel={trimestre},ylabel={flux versés à l'equity (M USD)},
  tick label style={font=\small},label style={font=\small},
  xmin=0.5,xmax=19.5,ymin=0,ymax=2.5,grid=major,grid style={gray!25},
  legend columns=3,legend style={font=\footnotesize,at={(0.5,-0.32)},anchor=north,draw=none,fill=none,/tikz/every even column/.append style={column sep=8pt}},
  table/col sep=comma]
\addplot[omDef,very thick,const plot mark mid,restrict x to domain=1:19] table[x=quarter,y=eq0]{figdata/markets-2/25-loans-clos-and-securitisation/paths.csv};
\addplot[omThm,very thick,const plot mark mid,restrict x to domain=1:19] table[x=quarter,y=eq1]{figdata/markets-2/25-loans-clos-and-securitisation/paths.csv};
\addplot[omProp,very thick,const plot mark mid,restrict x to domain=1:19] table[x=quarter,y=eq2]{figdata/markets-2/25-loans-clos-and-securitisation/paths.csv};
\legend{2~\% par an,4.8~\%,7~\%}
\end{axis}
\end{tikzpicture}

\omcaption{Le CLO du chapitre sous des taux de défaut annuels constants de 2~\%, 4.8~\% et 7~\%
(recouvrement de 70~\%, recouvrements réinvestis). En haut~: le ratio de surdimensionnement au niveau BB
face à son seuil~; une fois qu'il échoue, les détournements le maintiennent près du seuil aussi longtemps que
les intérêts le permettent. En bas~: la trésorerie trimestrielle des titres de capital, que les détournements
réduisent, à zéro dès le trimestre 12 pour 7~\% et à presque rien dès le trimestre 17 pour 4.8~\%~; la vente
des prêts du dernier trimestre n'est pas représentée. Illustratif~; données~: celles du
tutoriel du chapitre.}\label{fig:m2:loans-clos-and-securitisation:paths}
\end{omfigure}

\section{Qui détient quelle tranche}

Les obligations sont vendues à des investisseurs dont les besoins correspondent à leur risque~: les investisseurs dans les CLO américains
sont typiquement des banques, des fonds communs de placement, des compagnies d'assurance, des fonds de pension et des hedge funds.
Les tranches AAA rapportent un petit spread au-dessus du taux variable, en échange d'une protection
de 38~\% par subordination dans l'exemple du chapitre, ce qui convient aux investisseurs qui ont besoin de notations élevées
et d'actifs à taux variable~; les obligations mezzanine rapportent davantage pour une protection plus mince~; la
\omterm{def:m2:loans-clos-and-securitisation:tests}{tranche equity} est une position à effet de levier sur les prêts, d'environ dix fois, et perçoit le surplus
de spread, ce qui convient aux investisseurs capables de supporter de fortes pertes dans un mauvais cycle. Le gestionnaire
est rémunéré par des commissions qui passent avant toutes les obligations, quoi qu'il arrive à la \omterm{def:m2:loans-clos-and-securitisation:tests}{tranche equity}.

La structure rend les rendements de la \omterm{def:m2:loans-clos-and-securitisation:tests}{tranche equity} très sensibles aux défauts
(\cref{fig:m2:loans-clos-and-securitisation:irr}). Son taux de rendement interne tombe
de 19.5~\% sans défaut à 13.1~\% pour 2~\% par an, puis à 5.0~\% pour 4~\%, et il devient
négatif à partir d'environ 5~\%. Les tranches AAA, protégées par la subordination et par les
tests, qui transforment la trésorerie de la \omterm{def:m2:loans-clos-and-securitisation:tests}{tranche equity} en remboursements, sont remboursées intégralement dans tous les cas
étudiés par le tutoriel, même à 10~\% par an.

\begin{omfigure}
\begin{tikzpicture}
\begin{axis}[width=11cm,height=5cm,
  xlabel={taux de défaut constant (\% par an)},ylabel={TRI de la tranche equity (\%)},
  tick label style={font=\small},label style={font=\small},
  xmin=0,xmax=10,ymin=-40,ymax=25,grid=major,grid style={gray!25},
  axis y line*=left,
  legend columns=2,legend style={font=\footnotesize,at={(0.5,-0.3)},anchor=north,draw=none,fill=none,/tikz/every even column/.append style={column sep=8pt}},
  table/col sep=comma]
\addplot[omDef,very thick,mark=*,mark size=1.2pt] table[x=cdr,y=irr]{figdata/markets-2/25-loans-clos-and-securitisation/irr.csv};
\addplot[gray,thin,forget plot] coordinates {(0,0) (10,0)};
\addplot[omThm,thick,dashed] coordinates {(4.78,-40) (4.78,25)};
\legend{TRI de la tranche equity,tranche equity coupée en premier (4.78~\%)}
\end{axis}
\end{tikzpicture}

\omcaption{Taux de rendement interne de la \omterm{def:m2:loans-clos-and-securitisation:tests}{tranche equity} du CLO du chapitre, sur cinq ans avec
les prêts vendus au pair à la fin, en fonction d'un taux de défaut annuel constant (recouvrement de
70~\%). La ligne en pointillés marque le taux de défaut au-delà duquel les tests coupent à zéro la
trésorerie de la \omterm{def:m2:loans-clos-and-securitisation:tests}{tranche equity} pour un trimestre donné. Les rendements inférieurs à $-40\%$, à partir de 9~\% par an, sortent du
graphique. Illustratif~; données~: celles du tutoriel du chapitre.}\label{fig:m2:loans-clos-and-securitisation:irr}
\end{omfigure}

\section{Tutoriel~: une cascade de CLO face à une vague de défauts}\label{tut:m2:loans-clos-and-securitisation:waterfall}

\begin{tutorial}
\textbf{Objectif.} Construire la cascade, faire passer le CLO du chapitre par des taux de défaut
constants, observer le \omterm{def:m2:loans-clos-and-securitisation:tests}{test de surdimensionnement} détourner la trésorerie, et trouver le taux de défaut
à partir duquel la \omterm{def:m2:loans-clos-and-securitisation:tests}{tranche equity} ne reçoit plus de paiements.
\textbf{État final~:} \cref{fig:m2:loans-clos-and-securitisation:paths,fig:m2:loans-clos-and-securitisation:irr},
\cref{ex:m2:loans-clos-and-securitisation:base} et les chiffres du problème du week-end.
\begin{tutsteps}
\item \textbf{L'opération}~: obligations, déclencheurs et relevé trimestriel.
\omcode{code/firm/waterfall/firm_waterfall.py}{15}{44}{Obligations, conditions de l'opération et relevé d'une période.}\label{lst:m2:loans-clos-and-securitisation:deal}
\item \textbf{La cascade}~: défauts, réinvestissement, intérêts par catégorie, tests et
détournements (\texttt{pay\_down} rembourse les obligations dans l'ordre), la vente finale.
\omcode{code/firm/waterfall/firm_waterfall.py}{56}{94}{La cascade trimestrielle.}\label{lst:m2:loans-clos-and-securitisation:run}
\item \textbf{Les résultats}~: le taux de rendement interne de la \omterm{def:m2:loans-clos-and-securitisation:tests}{tranche equity} et le taux de défaut
le plus bas auquel elle est coupée, par dichotomie.
\omcode{code/firm/waterfall/firm_waterfall.py}{97}{117}{TRI de la tranche equity et taux de défaut de coupure.}\label{lst:m2:loans-clos-and-securitisation:irr}
\item \textbf{Exécuter} \texttt{clo\_demo.base\_case()}, \texttt{clo\_demo.problem()} et
\texttt{fig\_clo.py}.
\end{tutsteps}
\textbf{Pour aller plus loin.} Remplacer le taux de défaut constant par une vague, 2~\% par an
pendant deux ans, puis 10~\% pendant un an, puis 2~\%, et comparer le rendement de la \omterm{def:m2:loans-clos-and-securitisation:tests}{tranche equity} avec le
taux constant qui produit les mêmes défauts cumulés~; puis faire vendre les prêts à 95
au lieu du pair à la fin.
\end{tutorial}

\section{Construction~: le moteur de cascade du CLO}\label{bld:m2:loans-clos-and-securitisation:waterfall}

\begin{build}
\bfield{Rôle} Le desk de crédit de la firme miniature valorise les obligations et la \omterm{def:m2:loans-clos-and-securitisation:tests}{tranche equity} des CLO sous
des scénarios de défaut, et son équipe risque voit à quel point chaque opération est proche de ses déclencheurs.
\bfield{Interface} \texttt{Note(name, balance, spread, deferrable)}~; \texttt{Deal(notes, par, loan\_spread, rate, fee, oc, ic, recovery, quarters,
reinvest\_quarters)}~; \texttt{run(deal, cdr)} renvoyant un \texttt{Period} par
trimestre~; \texttt{pay\_down(notes, amount)}~; \texttt{equity\_irr(cash, invested)}~;
\texttt{cutoff\_cdr(deal)}.
\bfield{Règles} Trimestriel~; taux de défaut annuel constant~; recouvrements réinvestis au pair
pendant la période de réinvestissement, utilisés pour rembourser les obligations après celle-ci~; commissions passant avant toutes les obligations~;
intérêts différables de la tranche junior~; les rattrapages de surdimensionnement sont limités aux intérêts
disponibles~; un échec du \omterm{def:m2:loans-clos-and-securitisation:tests}{test de couverture des intérêts} détourne tous les intérêts restants~; prêts vendus au
pair à la fin.
\bfield{Tests d'acceptation} \texttt{code/firm/waterfall/tests/}~: sans défaut, chaque
catégorie est payée et la \omterm{def:m2:loans-clos-and-securitisation:tests}{tranche equity} reçoit le surplus de spread et son principal~; soldes
et paiements ne deviennent jamais négatifs~; au-dessus du taux de défaut de coupure, la \omterm{def:m2:loans-clos-and-securitisation:tests}{tranche equity} est coupée dans
un trimestre donné, en dessous, jamais.
\bfield{Pour aller plus loin} Pools au niveau des prêts avec notations et un seau CCC soumis à une décote dans les tests~;
scénarios de défaut et de remboursement anticipé issus d'un modèle de crédit~; cascade du principal et
critères de réinvestissement~; valorisation de chaque catégorie par simulation (livre~6).
\end{build}

\begin{omsources}
\item Federal Reserve Board, Financial Stability Report, mai 2026.
\item M.~Guse, W.~Park, Z.~Saravay et Y.~Yook, «~Collateralized loan obligations in the
Financial Accounts of the United States~», FEDS Notes, septembre 2019.
\item NAIC Capital Markets Bureau, «~Middle market collateralized loan obligations
primer~», mars 2025.
\item Chapman and Cutler, note client sur \emph{Loan Syndications and Trading
Association v.\ SEC} (D.C. Cir., 9 février 2018).
\end{omsources}

\section{Exercices}

\begin{exercise}[$\star$]\label{exo:m2:loans-clos-and-securitisation:1}
Calculer les quatre ratios de surdimensionnement du CLO du chapitre à l'émission.
\end{exercise}

\begin{exercise}[$\star$]\label{exo:m2:loans-clos-and-securitisation:2}
Pourquoi la charge d'intérêts d'un \omterm{def:m2:loans-clos-and-securitisation:loan}{prêt à effet de levier} augmente-t-elle lorsque la banque centrale relève ses taux, alors
que celle d'une obligation à taux fixe n'augmente pas~?
\end{exercise}

\begin{exercise}[$\star$]\label{exo:m2:loans-clos-and-securitisation:3}
Que retire un prêt \omterm{def:m2:loans-clos-and-securitisation:loan}{covenant-lite} à ses prêteurs~?
\end{exercise}

\begin{exercise}[$\star\star$]\label{exo:m2:loans-clos-and-securitisation:4}
Avec un \omterm{def:m2:corporate-bonds:covenant}{taux de recouvrement} de 70~\%, quelle part du portefeuille doit faire défaut avant
que le test de niveau A (seuil 1.18) ne soit en échec~?
\end{exercise}

\begin{exercise}[$\star\star$]\label{exo:m2:loans-clos-and-securitisation:5}
Pourquoi la cour a-t-elle jugé que la règle de rétention de risque ne convenait pas
aux gestionnaires de CLO sur le marché ouvert, et quel problème d'incitation la rétention de risque permet-elle de traiter~?
\end{exercise}

\begin{exercise}[$\star\star$]\label{exo:m2:loans-clos-and-securitisation:6}
L'equity a un effet de levier d'environ dix sur les prêts. En utilisant l'\cref{ex:m2:loans-clos-and-securitisation:base},
expliquer pourquoi son rendement monétaire vaut 18.2~\% alors que les prêts rapportent 3.5~\% au-dessus du taux variable.
\end{exercise}

\begin{exercise}[$\star\star\star$]\label{exo:m2:loans-clos-and-securitisation:7}
\emph{Programmation.} Mettre fin à la période de réinvestissement après deux ans au lieu de cinq, de sorte que
les recouvrements remboursent les billets, puis déterminer le nouveau taux de défaut limite. Expliquer le changement.
\end{exercise}

\begin{exercise}[$\star\star\star$]\label{exo:m2:loans-clos-and-securitisation:8}
\emph{Repérer l'erreur.} «~Les billets AAA d'un CLO sont aussi sûrs que les obligations AAA d'une
grande entreprise~: même notation, même risque.~» La corriger.
\end{exercise}

\section{Problème~: Le test qui déclenche}

\begin{problem}[{Problème du week-end --- quand l'equity cesse d'être versée}]\label{pb:m2:loans-clos-and-securitisation:1}
Un investisseur envisage l'equity du CLO du chapitre~: 500 millions de dollars de prêts à 3.5~\%
au-dessus d'un taux variable de 4~\%, des billets pour 450 millions de dollars répartis en cinq classes, 50 millions de dollars
d'equity, des frais senior de 0.45~\%, des seuils de surdimensionnement de 1.25, 1.18, 1.10 et
1.05 après les classes AA, A, BBB et BB, un seuil de couverture des intérêts de 1.20 après
la classe AA, des recouvrements de 70~\% réinvestis pendant cinq ans, et les prêts vendus au pair
à la fin.

\medskip\noindent\textbf{Partie I --- La structure.}
\begin{enumerate}
\item Que rapportent les prêts, et que coûtent les frais et les billets, chaque année~?
\item Que reçoit l'equity chaque année sans défaut, et quel
rendement monétaire cela représente-t-il~?
\item Quel est le rendement de l'equity sur cinq ans sans défaut~?
\item De combien le pair peut-il baisser avant que le test de niveau BB ne soit en échec~?
\item Quels défauts cumulés cela représente-t-il, avec un \omterm{def:m2:corporate-bonds:covenant}{taux de recouvrement} de 70~\%~?
\end{enumerate}

\medskip\noindent\textbf{Partie II --- Défauts.}
\begin{enumerate}[resume]
\item Quel est le rendement de l'equity à des taux de défaut constants de 2~\%, 4~\% et 6~\% par an~?
\item Quelle somme de trésorerie est détournée sur cinq ans à 4~\%, 6~\% et 10~\%~?
\item Déterminer le plus faible taux de défaut constant auquel l'equity ne reçoit rien
au cours d'un trimestre.
\item À ce taux, à quel trimestre l'equity est-elle coupée pour la première fois, et combien de pair
le portefeuille a-t-il perdu~?
\item Que se passe-t-il pour les billets BB à 10~\% par an~?
\end{enumerate}

\medskip\noindent\textbf{Partie III --- Leviers.}
\begin{enumerate}[resume]
\item Comment la fin de la période de réinvestissement après deux ans modifie-t-elle le seuil de coupure~?
\item Pourquoi les tests protègent-ils les billets AAA~?
\item Pourquoi un prêt qui recouvre 70~\% après défaut est-il moins dommageable pour les tests
qu'une obligation qui recouvre 40~\%~?
\item Que ferait un gestionnaire lorsque le test est près d'échouer, pour que l'equity
continue d'être versée, et quels sont les risques~?
\item Pourquoi est-il important pour l'equity que les prêts soient vendus au pair à la fin~?
\end{enumerate}

\medskip\noindent\textbf{Partie IV --- Jugement.}
\begin{enumerate}[resume]
\item À quels taux de défaut historiques comparer 4.78~\%, et pourquoi un taux constant
est-il optimiste en récession~?
\item Qui devrait détenir l'equity, et qui les billets AAA~?
\item Comment couvrir l'exposition de l'equity à une vague de défauts~?
\item Énoncer le \emph{résultat nommé}~: le taux de défaut auquel l'equity cesse de recevoir de la trésorerie,
et quand.
\item En une phrase~: à quoi sert le \omterm{def:m2:loans-clos-and-securitisation:tests}{test de surdimensionnement}~?
\end{enumerate}
\end{problem}

\section{Questions d'entretien}

\begin{interviewq}[$\star$ \iqroles{trader, chercheur}]\label{iq:m2:loans-clos-and-securitisation:1}
Qu'est-ce qu'un CLO, et qui détient ses tranches~?
\end{interviewq}

\begin{interviewq}[$\star$ \iqroles{chercheur, risque}]\label{iq:m2:loans-clos-and-securitisation:2}
Comment fonctionnent les \omterm{def:m2:loans-clos-and-securitisation:tests}{tests de surdimensionnement} et de couverture des intérêts~?
\end{interviewq}

\begin{interviewq}[$\star\star$ \iqroles{chercheur}]\label{iq:m2:loans-clos-and-securitisation:3}
Comment évalueriez-vous une \omterm{def:m2:loans-clos-and-securitisation:tests}{tranche equity} de CLO~?
\end{interviewq}

\begin{interviewq}[$\star\star$ \iqroles{trader, chercheur}]\label{iq:m2:loans-clos-and-securitisation:4}
Comparez une tranche de CLO à une tranche d'un \omterm{def:m2:credit-indices-and-tranches:index}{indice de crédit}. Qu'est-ce qui
distingue son risque~?
\end{interviewq}

\begin{interviewq}[$\star\star$ \iqroles{risque, banque}]\label{iq:m2:loans-clos-and-securitisation:5}
Quels sont les risques du marché des \omterm{def:m2:loans-clos-and-securitisation:loan}{prêts à effet de levier} pour la stabilité financière~?
\end{interviewq}

\begin{interviewq}[$\star\star\star$ \iqroles{développeur}]\label{iq:m2:loans-clos-and-securitisation:6}
Concevez un système qui exécute les \omterm{def:m2:loans-clos-and-securitisation:clo}{cascades des paiements} de 2\,000 CLO sous 10\,000
scénarios chacun, chaque nuit.
\end{interviewq}
